\section{Conclusão}
A biblioteca $\LL$ reúne segmentos afins racionais com endereço único e suporte diádico recuperável. A bijeção com jatos racionais produz montagens contínuas densas em $C(K)$ e uma representação ReLU rasa exata. Na aproximação uniforme de funções $C^2$, a racionalização preserva a constante $1/8$ do erro interpolatório.

O refinamento adaptativo combina majorantes intervalares da segunda derivada e racionalização nodal para obter erro até $\varepsilon_I+\varepsilon_R$, com término e número de folhas controlados sob as hipóteses declaradas. A contagem aritmética, a codificação e a seleção exata são desenvolvidas no segundo artigo; as agendas de altura e a teoria direta--inversa, no terceiro. Os quatro resultados de referência não reproduzidos nesses recortes têm texto e demonstração preservados no complemento do original.
